\documentclass[11pt,a4paper]{article}
\usepackage[T1]{fontenc}
\usepackage[utf8]{inputenc}
\usepackage{amsmath,amssymb,amsthm}
\usepackage[margin=1in]{geometry}
\usepackage[colorlinks=true,linkcolor=blue,urlcolor=blue]{hyperref}
\theoremstyle{plain}
\newtheorem{theorem}{Theorem}
\newtheorem{lemma}[theorem]{Lemma}
\newtheorem{proposition}[theorem]{Proposition}
\newtheorem{corollary}[theorem]{Corollary}
\theoremstyle{definition}
\newtheorem{remark}[theorem]{Remark}

\title{The empirical recurrence for OEIS A204777, proved by transfer matrix}
\author{Adrian Perez Fontelles\\ \small Independent researcher}
\date{7 September 2026}
\begin{document}
\maketitle

\begin{abstract}
OEIS A204777 counts $(n+2)\times5$ arrays over $\{0,\dots,2\}$ subject to a condition imposed on
every $3\times3$ subblock, and counted UP TO RELABELLING of the
alphabet; it carries an empirical recurrence of order $64$ contributed by
R.~H.~Hardin. It is true, and it is decidable rather than empirical. What is counted is
equality patterns, which is not a walk count --- but it is a fixed rational combination of
walk counts, one for each alphabet size, and that combination is again $C$-finite. Whether the
conjectured recurrence annihilates it is then decided by a finite exact computation, with the
Cayley--Hamilton theorem bounding the work.
\end{abstract}

\noindent\small 2020 Mathematics Subject Classification. 05A15, 05B45, 15B36.\normalsize

\section{The conjecture}

OEIS A204777 is ``Number of (n+2)X5 0..2 arrays with no 3X3 subblock having three equal diagonal elements or three equal antidiagonal elements, and new values 0..2 introduced in row major order.''. It has offset $1$ and begins
\[
1180224, 160435446, 21570701058, 2890032532056, 387664140116880, 51989763458957400, 6972630432052561776, 935137622698422608736,\ \dots
\]
The entry states, as a conjecture contributed by its author and never marked settled:
\begin{quote}
Empirical: a(n) = 139*a(n-1) -91120*a(n-3) +133552*a(n-4) +42248076*a(n-5) -149986356*a(n-6) -7637679060*a(n-7) +313732968864*a(n-8) -38217278758320*a(n-9) +156983430294144*a(n-10) +3146043316395312*a(n-11) -22524502662024800*a(n-12) +777226283328788768*a(n-13) -2923856345749409280*a(n-14) -55967658209930232128*a(n-15) +288975848517914503808*a(n-16) +38954422078319864832*a(n-17) -6182801204412109345536*a(n-18) +16684100050565913346560*a(n-19) +250356441942356912185344*a(n-20) -2383818791297995227921408*a(n-21) +1201231326963892811839488*a(n-22) +86802835054091019154550784*a(n-23) -184887655399216097211469824*a(n-24) +471923427737198981007802368*a(n-25) -1105116310670533761772781568*a(n-26) +8839932835611664798842028032*a(n-27) -144669014392674806035633668096*a(n-28) +450994890935237153381157961728*a(n-29) -291707817095511715532134809600*a(n-30) -27193350116155962634997969977344*a(n-31) +14610495860844431841467075395584*a(n-32) +71611394407979021390497695399936*a(n-33) -39885557889095088580277437464576*a(n-34) -505135228048555321215662325497856*a(n-35) +630569166822546215503538014912512*a(n-36) -248452930839886449596900988420096*a(n-37) +526298018115099117274747744813056*a(n-38) +61968287910240254616686971481751552*a(n-39) -31166308186060813617619771053834240*a(n-40) -119157499688627783035052503225860096*a(n-41) +68879591620458111760382530726920192*a(n-42) +1303817841279799249962996202790191104*a(n-43) -991814099212870169326345207098114048*a(n-44) -1676724409254838964837194313056124928*a(n-45) +790061313693482027810811270301483008*a(n-46) -39260160904271413856472704669454434304*a(n-47) +23310288012876239389061106505497968640*a(n-48) +82112809486892283035692868064126173184*a(n-49) -47379038370425292034007364935784333312*a(n-50) -147039553387406857205575820606240194560*a(n-51) +95223216026612306973536229244853551104*a(n-52) -81546391794155061437501143765217181696*a(n-53) +86118064909915539137280691409145298944*a(n-54) +2310263373801310827876943006321952686080*a(n-55) -1357420446491639344646358413831609253888*a(n-56) -67278639720595731178377088196319117312*a(n-57) -154291462019127049141482660348479668224*a(n-58) -3385126747711336339558048425047293427712*a(n-59) +1821392370949169125647162467129896206336*a(n-60) +104474909490376310181853533272788697088*a(n-61) +812752711529833869502920372858365411328*a(n-63) -425467862411590884706226772368808738816*a(n-64)
\end{quote}
As of the ``Last modified'' line on the live entry (Oct 03 2025 11:30:01, revision 9) this is still
recorded as empirical, and nothing on the entry records it as proved.

\section{Patterns are a combination of counts}

The clause ``new values $0..2$ introduced in row major order'' says that the arrays
counted are the canonical representatives of their EQUALITY PATTERNS: two arrays that differ
only by a permutation of the alphabet are one object. Write $N_j$ for the number of admissible
patterns using exactly $j$ letters and $L_i$ for the number of admissible ARRAYS over an
alphabet of $i$ letters. An array over $i$ letters is a pattern together with an injection of
its letters into the alphabet, so
\[
L_i\;=\;\sum_{j}N_j\,i(i-1)\cdots(i-j+1),
\]
a triangular system with nonzero diagonal, hence invertible over $\mathbb{Q}$. The entry
counts $a=N_1+\dots+N_{3}$, so
\[
a\;=\;\bigl(\mathbf{1}^{\!\top}F^{-1}\bigr)L ,
\]
a FIXED rational combination of the $L_i$, where $F$ is the matrix of falling factorials.
Each $L_i$ is a walk count on its own transfer graph, so $a$ is a walk count on the disjoint
union of those graphs with the combination's coefficients placed in the starting vector; the
coefficients are rational and some are negative, which costs nothing in what follows.

This step needs the condition itself to be unchanged by permuting the alphabet, and it is:
the condition below is stated through equality between entries and never names a particular
value or compares sizes.

\section{Three consecutive lines}

The entry's condition constrains every $3\times3$ subblock, which spans three consecutive
rows and three consecutive columns. Writing an array as a sequence of rows
$u_1,u_2,\dots$ over $\{0,\dots,2\}^{5}$, the condition
\[
\text{not }\bigl(\text{the diagonal or the antidiagonal of }g\text{ is constant}\bigr)
\]
involves only three consecutive rows. Take as vertices the ordered PAIRS $(r,s)$ of
rows and put an edge $(r,s)\to(s,t)$ exactly when the condition holds for the triple
$(r,s,t)$ at every admissible column; an array with $L$ rows is a walk of length $L-2$.
Doing that for each alphabet size $1,\dots,3$ and combining as above gives
\[
a(n)\;=\;\frac{1}{1}\,\iota^{\!\top}M^{\,n}\tau
\]
on $S=59050$ vertices in total, the exponent fixed by the entry's own shape and checked against
its published terms. States lying on no walk from $\iota$ to $\tau$ are removed first; that
changes no value and only shrinks the computation.

\section{The proof}

\begin{lemma}\label{lem:crit}
Let $q(t)=t^{64}-\sum_i c_it^{64-i}$ be the characteristic polynomial of the
conjectured recurrence. Then the residual $a(n)-\sum_ic_ia(n-i)$ equals
$\iota^{\!\top}M^{\,j}q(M)\tau$ for the corresponding walk index $j$, so the
recurrence holds from some point on if and only if $\iota^{\!\top}M^{j}q(M)\tau=0$ for all
large $j$; and it is enough to examine $j<S$.
\end{lemma}

\begin{proof}
Factor $M^{j}$ out of $M^{j+64}-\sum_ic_iM^{j+64-i}$. For the bound, $M^{S}$
is an integer combination of $I,M,\dots,M^{S-1}$ by Cayley--Hamilton, so the numbers
$u_j=\iota^{\!\top}M^{j}q(M)\tau$ satisfy the monic linear recurrence given by the
characteristic polynomial of $M$; $S$ consecutive zeros force every later one, and the last
nonzero $u_j$ pins the threshold exactly.
\end{proof}

\begin{theorem}
\[
a(n)\;=\;139\,a(n-1) - 91120\,a(n-3) + 133552\,a(n-4) + 42248076\,a(n-5) - 149986356\,a(n-6) - 7637679060\,a(n-7) + 313732968864\,a(n-8) - 38217278758320\,a(n-9) + 156983430294144\,a(n-10) + 3146043316395312\,a(n-11) - 22524502662024800\,a(n-12) + 777226283328788768\,a(n-13) - 2923856345749409280\,a(n-14) - 55967658209930232128\,a(n-15) + 288975848517914503808\,a(n-16) + 38954422078319864832\,a(n-17) - 6182801204412109345536\,a(n-18) + 16684100050565913346560\,a(n-19) + 250356441942356912185344\,a(n-20) - 2383818791297995227921408\,a(n-21) + 1201231326963892811839488\,a(n-22) + 86802835054091019154550784\,a(n-23) - 184887655399216097211469824\,a(n-24) + 471923427737198981007802368\,a(n-25) - 1105116310670533761772781568\,a(n-26) + 8839932835611664798842028032\,a(n-27) - 144669014392674806035633668096\,a(n-28) + 450994890935237153381157961728\,a(n-29) - 291707817095511715532134809600\,a(n-30) - 27193350116155962634997969977344\,a(n-31) + 14610495860844431841467075395584\,a(n-32) + 71611394407979021390497695399936\,a(n-33) - 39885557889095088580277437464576\,a(n-34) - 505135228048555321215662325497856\,a(n-35) + 630569166822546215503538014912512\,a(n-36) - 248452930839886449596900988420096\,a(n-37) + 526298018115099117274747744813056\,a(n-38) + 61968287910240254616686971481751552\,a(n-39) - 31166308186060813617619771053834240\,a(n-40) - 119157499688627783035052503225860096\,a(n-41) + 68879591620458111760382530726920192\,a(n-42) + 1303817841279799249962996202790191104\,a(n-43) - 991814099212870169326345207098114048\,a(n-44) - 1676724409254838964837194313056124928\,a(n-45) + 790061313693482027810811270301483008\,a(n-46) - 39260160904271413856472704669454434304\,a(n-47) + 23310288012876239389061106505497968640\,a(n-48) + 82112809486892283035692868064126173184\,a(n-49) - 47379038370425292034007364935784333312\,a(n-50) - 147039553387406857205575820606240194560\,a(n-51) + 95223216026612306973536229244853551104\,a(n-52) - 81546391794155061437501143765217181696\,a(n-53) + 86118064909915539137280691409145298944\,a(n-54) + 2310263373801310827876943006321952686080\,a(n-55) - 1357420446491639344646358413831609253888\,a(n-56) - 67278639720595731178377088196319117312\,a(n-57) - 154291462019127049141482660348479668224\,a(n-58) - 3385126747711336339558048425047293427712\,a(n-59) + 1821392370949169125647162467129896206336\,a(n-60) + 104474909490376310181853533272788697088\,a(n-61) + 812752711529833869502920372858365411328\,a(n-63) - 425467862411590884706226772368808738816\,a(n-64)
\]
for every $n>64$.
\end{theorem}

\begin{proof}
The vector $w=q(M)\tau$ was formed by $64$ matrix--vector products in exact integer
arithmetic and $\iota^{\!\top}M^{j}w$ evaluated for $j=0,1,\dots$, again exactly. Those
integers vanish from the index corresponding to $n=64+1$ onwards, and the run was
continued until the working vector was identically zero, which settles all larger $j$ at once.
By Lemma~\ref{lem:crit} the recurrence holds for every $n>64$.
\end{proof}

\section{Verification}

Three checks, each independent of the algebra above.

First, the transfer model reproduces all $10$ terms the entry publishes, exactly, in
integer arithmetic. That check earned its place here. The neighbour offsets are named in the
array's own frame --- horizontal means along a row --- but when the entry fixes the first
dimension, as in ``$2\times n$'', the walk runs along columns and the two components of every
offset swap. Sets such as king-move, or horizontal together with vertical, are unchanged by
that swap and hide the error completely; a set such as horizontal, diagonal and antidiagonal
is not, and the counts came out as those of a different member of the same family until the
offsets were transposed.

Second, the conjectured recurrence was evaluated directly on the entry's own published terms,
with no matrices involved, and holds wherever the proved range and the data overlap.

Third, the annihilation test of Lemma~\ref{lem:crit} was carried out in exact integer
arithmetic on the whole vector rather than on a sampled prefix.

\begin{thebibliography}{9}
\bibitem{oeis} The OEIS Foundation, \emph{The On-Line Encyclopedia of Integer
Sequences}, \url{https://oeis.org}, 2026. Sequence A204777.
\bibitem{stanley} R.~P.~Stanley, \emph{Enumerative Combinatorics}, Volume 1, 2nd ed.,
Cambridge University Press, 2011. (Transfer-matrix method, Section 4.7.)
\bibitem{fs} P.~Flajolet and R.~Sedgewick, \emph{Analytic Combinatorics}, Cambridge
University Press, 2009.
\bibitem{horn} R.~A.~Horn and C.~R.~Johnson, \emph{Matrix Analysis}, 2nd ed., Cambridge
University Press, 2013.
\end{thebibliography}

\end{document}
